\section{Kết luận}

Qua bài lab này, tôi đã trải qua toàn bộ chu trình cài đặt, cấu hình và vận hành
một cụm Hadoop — từ chế độ pseudo-distributed đơn giản trên máy cá nhân cho đến
mô hình fully-distributed nhiều node. Đây là nền tảng thực hành không thể thiếu
để hiểu sâu hơn về kiến trúc của hệ sinh thái Big Data.

\textbf{Những kiến thức thu được:}

\begin{itemize}
  \item \textbf{Kiến trúc HDFS}: tôi hiểu rõ vai trò phân công giữa NameNode (quản
        lý metadata và namespace) và DataNode (lưu trữ dữ liệu thực tế dạng block),
        cũng như cơ chế replication giúp đảm bảo tính chịu lỗi (fault tolerance).
        Nguyên tắc \textit{write once, read many} của HDFS phù hợp với các workload
        phân tích dữ liệu lớn.

  \item \textbf{Hệ thống YARN}: tôi nắm được mô hình quản lý tài nguyên hai tầng
        của YARN — ResourceManager toàn cục điều phối tài nguyên, còn ApplicationMaster
        từng job tự quản lý lifecycle của các task. Kiến trúc này tách biệt quản lý
        tài nguyên khỏi quản lý job, giúp YARN hỗ trợ nhiều framework (MapReduce,
        Spark, Flink) trên cùng một cluster.

  \item \textbf{Mô hình MapReduce}: Qua bài toán WordCount, tôi thực hành
        triển khai mô hình lập trình phân tán theo paradigm Map-Reduce. Tôi hiểu rõ 
        cách dữ liệu được phân mảnh (Split) và xử lý độc lập ở giai đoạn Map, sau đó
        được nhóm (Shuffle \& Sort) trước khi tổng hợp ở giai đoạn Reduce. Việc tùy biến 
        Custom Comparator cũng giúp tôi hiểu thêm luồng xử lý sâu bên trong Hadoop.

  \item \textbf{Quản lý phân quyền HDFS}: hệ thống phân quyền của HDFS tương tự
        Unix (owner, group, others) nhưng vận hành trong môi trường phân tán, đòi
        hỏi sự nhất quán về tên người dùng trên nhiều node.
\end{itemize}

\textbf{Khó khăn và bài học kinh nghiệm:}

Khó khăn lớn nhất tôi gặp phải là sự không tương thích giữa Cluster ID sau khi
format lại NameNode và ID cũ còn lưu trong thư mục DataNode, khiến DataNode từ
chối kết nối vào cluster. Bài học rút ra là: trong môi trường production, không bao
giờ format lại NameNode trừ khi xóa sạch toàn bộ dữ liệu DataNode đồng thời. Ngoài
ra, việc sử dụng Arch Linux — hệ điều hành không phổ biến trong cộng đồng Hadoop —
đôi khi gây ra các vấn đề về tên giao diện mạng không theo chuẩn truyền thống, đòi
hỏi phải tìm hiểu thêm để cấu hình đúng.

Kinh nghiệm thực hành qua bài lab này là bước đệm quan trọng để tiếp cận các chủ
đề nâng cao hơn trong môn học, bao gồm tối ưu hóa job MapReduce, tuning cấu hình YARN,
và tích hợp Hadoop với các hệ thống lưu trữ và xử lý hiện đại như Hive, HBase, và
Kafka.
